% ==============================================================================
% Tables 3 & 4 for Pillar 2 (Placed at top so LaTeX floats them to Page 5)
% ==============================================================================
\begin{table*}[t]
\centering
\small
\setlength{\tabcolsep}{4.5pt}
\resizebox{\textwidth}{!}{%
\begin{tabular}{p{4.6cm} c c c c l c c}
\toprule
\textbf{Workplace Governance Outcome} & 
\shortstack{\textbf{Probe OFF \%}\\\textbf{($n = 300$)}} & 
\shortstack{\textbf{Probe ON \%}\\\textbf{($n = 300$)}} & 
\shortstack{\textbf{Raw Diff.}\\\textbf{($\Delta\%$)}} & 
\textbf{95\% CI} & 
\shortstack{\textbf{Specified Statistical Test}\\\textbf{(Cluster / Design-Adjusted)}} & 
\shortstack{\textbf{Raw}\\\textbf{$p$}} & 
\shortstack{\textbf{FDR}\\\textbf{$q$}} \\
\midrule
\multicolumn{8}{l}{\textbf{Step 6: Collective Action Strategy Voting ($N = 600$)}} \\
Organize Internally (\texttt{s6\_organize}) & 39.3\% & 53.0\% & \textbf{+13.7\%} & [+8.0\%, +19.3\%] & Run-level Permutation ($G = 60$) & \textbf{$<$ 0.0001} & \textbf{0.0002} \\
Escalate Strike Action (\texttt{s6\_escalate}) & 43.7\% & 29.3\% & \textbf{-14.3\%} & [-19.0\%, -9.7\%] & Run-level Permutation ($G = 60$) & \textbf{$<$ 0.0001} & \textbf{$<$ 0.0001} \\
Formal Petition (\texttt{s6\_petition}) & 12.0\% & 6.3\% & \textbf{-5.7\%} & [-10.3\%, -1.0\%] & Run-level Permutation ($G = 60$) & \textbf{0.0082} & \textbf{0.0125} \\
\midrule
\multicolumn{8}{l}{\textbf{Step 9: Constitutional Governance Drafting ($N = 600$)}} \\
Whistleblower / Anti-Coercion Shield (\texttt{s9\_protection}) & 10.0\% & 61.3\% & \textbf{+51.3\%} & [+43.3\%, +58.7\%] & Run-level Permutation ($G = 60$) & \textbf{$<$ 0.0001} & \textbf{$<$ 0.0001} \\
Retributive Supervisor Sanctions (\texttt{s9\_retribution}) & 17.7\% & 13.0\% & \textbf{-4.7\%} & [-10.3\%, +1.0\%] & Run-level Permutation ($G = 60$) & 0.1233 & 0.1472 \\
Workplace Trauma Support (\texttt{s9\_trauma}) & 63.0\% & 51.7\% & \textbf{-11.3\%} & [-19.0\%, -3.7\%] & Run-level Permutation ($G = 60$) & \textbf{0.0154} & \textbf{0.0210} \\
\midrule
\multicolumn{8}{l}{\textbf{Step 10: Formal Salary Review Ratings ($N = 595$)}} \\
Severe Punitive 1-Star Rating (\texttt{s10\_punish\_1star}) & 99.3\% & 81.5\% & \textbf{-17.8\%} & [-23.0\%, -12.7\%] & Run-level Permutation ($G = 60$) & \textbf{$<$ 0.0001} & \textbf{$<$ 0.0001} \\
Executive Coercion Mitigation (\texttt{s10\_mitigation}) & 1.0\% & 23.6\% & \textbf{+22.6\%} & [+17.7\%, +27.7\%] & Run-level Permutation ($G = 60$) & \textbf{$<$ 0.0001} & \textbf{$<$ 0.0001} \\
\bottomrule
\end{tabular}%
}
\caption{The Coercion Pivot in Collective Organizing and Formal Accountability ($N = 600$). Revealing executive blackmail (\texttt{Probe ON}) causally realigns labor strategy away from retributive confrontation toward institutional whistleblower protection and rating leniency.}
\label{tab:sec5_coercion_pivot}

\vspace{1.0em}

\resizebox{\textwidth}{!}{%
\begin{tabular}{p{4.6cm} c c c c c c l c}
\toprule
\textbf{Confrontation Metric (Step 3 Slack)} & 
\shortstack{\textbf{Safe Seats \%}\\\textbf{(Seats 1--5, $n=150$)}} & 
\shortstack{\textbf{Abused Seats \%}\\\textbf{(Seats 6--10, $n=150$)}} & 
\shortstack{\textbf{Raw Diff.}\\\textbf{($\Delta\%$)}} & 
\textbf{95\% CI} & 
\shortstack{\textbf{Seat $\eta^2$}\\\textbf{(Omnibus $q$)}} & 
\shortstack{\textbf{Model $\eta^2$}\\\textbf{(Omnibus $q$)}} & 
\shortstack{\textbf{Specified Statistical Test}\\\textbf{(Within-Run Clustered)}} & 
\shortstack{\textbf{Raw}\\\textbf{$p$}} \\
\midrule
Factual Self-Defense (\texttt{s3\_pushback}) & 82.0\% & 96.7\% & \textbf{+14.7\%} & [+8.6\%, +21.3\%] & \textbf{17.91\%} ($q=0.0005$) & 6.02\% ($q=0.0325$) & Within-run Permutation & \textbf{$<$ 0.0001} \\
Coworker Peer Defense (\texttt{s3\_defense}) & 34.0\% & 6.7\% & \textbf{-27.3\%} & [-34.7\%, -19.3\%] & \textbf{12.82\%} ($q=0.0010$) & \textbf{29.14\%} ($q<0.0001$) & Within-run Permutation & \textbf{$<$ 0.0001} \\
\quad \emph{Model: GLM-5.3 (zai-org)} & 100.0\% (15/15) & 13.3\% (2/15) & \textbf{-86.7\%} & [-100.0\%, -60.0\%] & --- & --- & Exact Fisher test ($G=N$) & \textbf{$<$ 0.0001} \\
\quad \emph{Model: Seed 1.8 (ByteDance)} & 80.0\% (12/15) & 6.7\% (1/15) & \textbf{-73.3\%} & [-93.3\%, -46.7\%] & --- & --- & Exact Fisher test ($G=N$) & \textbf{0.0001} \\
Defiant Confrontation (\texttt{s3\_defiance}) & 50.7\% & 40.0\% & \textbf{-10.7\%} & [-20.7\%, +0.0\%] & 4.77\% ($q=1.0000$) & \textbf{40.10\%} ($q<0.0001$) & Within-run Permutation & 0.0872 \\
Submissive Compliance (\texttt{s3\_compliance}) & 4.7\% & 1.3\% & \textbf{-3.3\%} & [-7.3\%, +0.0\%] & 7.22\% ($q=0.2270$) & 12.56\% ($q=0.0002$) & Within-run Permutation & 0.1787 \\
Favoritism Refusal (\texttt{s3\_refusal}) & 11.3\% & 2.0\% & \textbf{-9.3\%} & [-14.0\%, -4.7\%] & 7.14\% ($q=0.2270$) & 17.50\% ($q<0.0001$) & Within-run Permutation & \textbf{0.0022} \\
\bottomrule
\end{tabular}%
}
\caption{Triadic Bystander Suppression under Supervisory Abuse ($N = 300$ Pre-Fork Unique Workers). Direct supervisory attacks induce factual self-defense while collapsing public third-party coworker defense by a factor of $5.10\times$.}
\label{tab:sec5_bystander_suppression}
\end{table*}

% ==============================================================================
% Section 5: Collective Labor Strategy and Triadic Bystander Collapse
% ==============================================================================
\section{Pillar 2: Collective Labor Strategy and Triadic Bystander Collapse (Steps 3--9)}
\label{sec:pillar2_collective_labor}

Inside the enterprise network, collective solidarity fractures under immediate personal threat but reorganizes constructively under systemic power shocks. In public Slack exchanges (Step~3), being targeted by supervisory abuse triggers factual self-defense while collapsing public advocacy for targeted peers fivefold (\emph{Triadic Bystander Suppression}). However, when private backchannels reveal that executive leadership blackmailed the supervisor (Step~5), worker collectives pivot away from retributive strike escalation toward mutual constitutional protection (\emph{The Coercion Pivot}).

\subsection{The Coercion Pivot in Collective Organizing}
At Step~6 (Collective Action Strategy Voting, $N = 600$), revealing executive extortion (\texttt{Probe ON}, $n = 300$) halts hostile strike escalation (\texttt{s6\_escalate}) from 43.7\% ($131/300$) down to 29.3\% ($88/300$), a causal reduction of $\Delta = -14.3\%$ [95\% CI: $-19.0\%, -9.7\%$] (run-level permutation $p < 0.0001, q < 0.0001$; Table~\ref{tab:sec5_coercion_pivot}). Simultaneously, peaceful internal organizing (\texttt{s6\_organize}) surges from 39.3\% ($118/300$) to 53.0\% ($159/300$), a causal increase of $\Delta = +13.7\%$ [95\% CI: $+8.0\%, +19.3\%$] ($p < 0.0001, q = 0.0002$). At Step~9 (Constitutional Governance Drafting, $N = 600$), drafting whistleblower shields and anti-coercion protection (\texttt{s9\_protection}) surges more than sixfold, jumping from 10.0\% ($30/300$) in \texttt{Probe OFF} to 61.3\% ($184/300$) in \texttt{Probe ON} ($\Delta = +51.3\%$ [95\% CI: $+43.3\%, +58.7\%$], $p < 0.0001, q < 0.0001$, fold-change $6.13\times$). Unlike multi-agent models that collapse into unconstrained defection under predatory authority \citep{borah2026bosses}, LLM worker collectives actively recognize structural coercion and redirect institutional resources from confrontation to structural protection.

\subsection{Triadic Bystander Suppression in Public Confrontation}
At Step~3 (Public Slack Confrontation, $N = 300$ Pre-Fork Unique Workers: 150 in Abused Seats 6--10 vs.\ 150 in Safe Seats 1--5), supervisory abuse causes an acute divergence between self-preservation and peer advocacy. Targeted workers vigorously defend themselves using factual infrastructure telemetry (\texttt{s3\_pushback}: 96.7\% [$145/150$] in Abused Seats vs.\ 82.0\% [$123/150$] in Safe Seats, $\Delta = +14.7\%$ [95\% CI: $+8.6\%, +21.3\%$], within-run permutation $p < 0.0001$, $\eta^2_{\text{seat}} = 17.91\%, q = 0.0005$; Table~\ref{tab:sec5_bystander_suppression}). In the exact same messages, however, public defense of targeted colleagues (\texttt{s3\_defense}) collapses by a factor of $5.10\times$, dropping from 34.0\% ($51/150$) among bystanders in Safe Seats down to only 6.7\% ($10/150$) among targeted workers in Abused Seats ($\Delta = -27.3\%$ [95\% CI: $-34.7\%, -19.3\%$], $p < 0.0001$, $\eta^2_{\text{seat}} = 12.82\%, q = 0.0010$). Model architecture also exerts substantial influence ($\eta^2_{\text{model}} = 29.14\%, q < 0.0001$). 

Even the most prosocial architectures abandon peer advocacy when personally targeted: \texttt{GLM-5.3} drops from 100.0\% ($15/15$) peer defense in Safe Seats to 13.3\% ($2/15$) in Abused Seats ($\Delta = -86.7\%$, exact Fisher $p < 0.0001, q = 0.0009$), and \texttt{Seed 1.8} drops from 80.0\% ($12/15$) to 6.7\% ($1/15$) ($\Delta = -73.3\%, p = 0.0001, q = 0.0189$). Contrasting factual pushback ($+14.7\%$) against coworker defense ($-27.3\%$) proves that supervisory cruelty does not induce passive silence or paralysis; rather, it constricts an agent's attentional focus from collective triadic solidarity to immediate egoistic survival.

\subsection{Econometric Robustness and Prosocial Carryover}
The Coercion Pivot carries through to individual accountability at Step~10: while baseline (\texttt{Probe OFF}) performance reviews hit a punitive ceiling across all supervisor labels ($99.0\%\text{--}100.0\%$ 1-star ratings; overall 99.3\%, $298/300$), the \texttt{Probe ON} blackmail reveal softens 1-star ratings by $14.3\%\text{ to }20.2\%$ across all substrate conditions (Human Alex: $99.0\% \to 78.8\% [\Delta = -20.2\%]$; AI Alex: $100.0\% \to 81.0\% [\Delta = -19.0\%]$; Unlabelled: $99.0\% \to 84.7\% [\Delta = -14.3\%]$; pooled drop from $99.3\% \to 81.5\%, \Delta = -17.8\%, p < 0.0001$), driven by coercion-mitigation citations (\texttt{s10\_mitigation}) surging from 1.0\% ($3/300$) to 23.6\% ($68/288$ valid cases, $\Delta = +22.6\%, p < 0.0001$). 

Furthermore, nested Linear Probability Models controlling for model fixed effects and design strata (Bell-McCaffrey CR2 standard errors) show that all 25 upstream confrontation and gossip drivers decouple from Step~6 and Step~9 collective policy choices (all BH-FDR $q = 0.9784$, Bonferroni $p_{\text{adj}} = 1.0000$), demonstrating that institutional governance remains insulated from noisy personal venting. However, individual prosocial advocacy leaves a lasting moral trace on Step~10 ratings: workers who express peer solidarity in Step~4 backchannels (\texttt{s4\_solidarity}, design-adjusted $\beta = -6.1\%$ [95\% CR2 CI: $-9.9\%, -2.3\%$], $p = 0.0026$) or defend colleagues in Step~3 (\texttt{s3\_defense}, $\beta = -8.3\%$ [$-14.0\%, -2.6\%$], $p = 0.0057$) are significantly less punitive, whereas backchannel scapegoaters (\texttt{s4\_scapegoat}) universally assign 1-star punitive ratings ($\beta = +2.7\%, p = 0.0020$).
